
% Codifica e Lingua
\usepackage[utf8]{inputenc}
\usepackage[main=italian]{babel}

% Margini (adatti per la rilegatura)
\usepackage[top=2.5cm, bottom=2.5cm, left=3cm, right=2.5cm]{geometry}

% Spaziatura tra i paragrafi (stile moderno)
\usepackage{parskip}

% Gestione Immagini
\usepackage{graphicx}
\graphicspath{{immagini/}} % Dice a LaTeX di cercare le foto direttamente qui

% Gestione del codice sorgente
\usepackage{listings}
\usepackage{xcolor}

% csquotes è fortemente consigliato insieme a babel e biblatex
\usepackage{csquotes}

% Bibliografia (moderna con Biber)
\usepackage[backend=biber, style=numeric, sorting=nty]{biblatex}

% Link cliccabili nel PDF (caricalo sempre per ultimo)
\usepackage[colorlinks=true, allcolors=black]{hyperref}

\usepackage{booktabs}
\usepackage{tabularx}
\usepackage{float}


\lstset{
    language=Python,
    basicstyle=\ttfamily\small,
    keywordstyle=\color{blue},
    stringstyle=\color{red},
    commentstyle=\color{green!60!black},
    breaklines=true,
    aboveskip=\smallskipamount,
    belowskip=\smallskipamount
}